\chapter{Notation and the Rosetta Stone}
\label{app:notation}

This book fixes one notation across six research papers and four literatures that did
not share one. This appendix records the choices and translates them.

\section{Notation used in this book}

\begin{center}
\footnotesize
\begin{tabular}{@{}lll@{}}
\toprule
Symbol & Meaning & First used \\
\midrule
$G$, $H$ & the quadratic form (Hessian); $G$ in Parts I--II, $H$ in Part III
         & Ch.~\ref{ch:oneproblem}, Ch.~\ref{ch:parametric} \\
$a$, $d$ & the linear term; $-a\T x$ in Parts I--II, $-\lambda d\T x$ in Part III
         & Ch.~\ref{ch:oneproblem}, Ch.~\ref{ch:parametric} \\
$C$, $b$ & constraint normals (columns) and right-hand sides, $C\T x \ge b$
         & Ch.~\ref{ch:oneproblem} \\
$m_{\mathrm{eq}}$ & number of leading constraints treated as equalities
         & Ch.~\ref{ch:oneproblem} \\
$A$ & in Parts I--II, the working-set normals $C_{\cdot\Active}$;
      in Part III, the equality system & Ch.~\ref{ch:oneproblem} \\
$\Active$ & working set (algorithm state) or active set (property of a point)
         & Def.~\ref{def:activeset} \\
$\Free$, $\Bound$ & free and blocked (bound) index sets
         & Ch.~\ref{ch:geometry} \\
$\lambda$, $u$, $\nu$, $\zeta$ & multipliers; $\lambda$ full-length, $u$ working-set,
      $\nu,\zeta$ in Part III & Ch.~\ref{ch:certificate} \\
$x_u = G^{-1}a$ & the unconstrained minimiser & Prop.~\ref{prop:sub} \\
$A\T G^{-1}A$ & the dual Hessian of a working set & Def.~\ref{def:dualhessian} \\
$J$, $R$ & the carried factors, $J\T GJ = I$, $J\T A = [R;0]$
         & Ch.~\ref{ch:walking} \\
$z$, $r$ & the primal and dual step directions & \eqref{eq:dzr} \\
$s$ & slacks (Part II) or the reduced gradient $Ax - b$ (bound-constrained case)
         & Ch.~\ref{ch:walking}, Ch.~\ref{ch:geometry} \\
$\kappa$ & spectral condition number $\lambda_{\max}/\lambda_{\min}$
         & Ch.~\ref{ch:freeblock} \\
$\Sig$, $\mu$, $w$ & covariance, expected returns, portfolio weights
         & Ch.~\ref{ch:oneproblem} \\
$X$, $y$, $\beta$ & design, response, regression coefficients
         & Ch.~\ref{ch:homotopy} \\
\bottomrule
\end{tabular}
\end{center}

Two conventions inherited from the reference implementations and used throughout
Parts~I and~II: the linear term is \emph{subtracted}, so the unconstrained minimiser is
$G^{-1}a$ rather than $-G^{-1}a$; and constraints are held \emph{column-wise}, so the
system is $C\T x \ge b$ rather than $Cx \le b$.

One genuine collision the reader should watch for. In Parts~I and~II, $A$ is the matrix
of working-set normals. In Part~III, following the parametric literature, $A$ is the
equality-constraint system and the working set is carried by the partition
$(\Free,\Bound,\mathcal{S})$. The chapters say which is in force; the collision is
retained rather than resolved because both usages are standard in their own literature
and a reader going on to the primary sources will meet both.

\section{The Rosetta stone}

Each row is one component of the common scheme of Chapter~\ref{ch:parametric}; the
columns give the name it goes by in the four literatures that developed it
independently. The leftmost column is the neutral optimization term of which all four
are special cases.
\index{Rosetta stone}

\begin{center}
\footnotesize
\begin{tabular}{@{}p{0.15\textwidth}p{0.19\textwidth}p{0.19\textwidth}p{0.17\textwidth}p{0.19\textwidth}@{}}
\toprule
Common scheme & Portfolio selection (CLA) & Regression (LASSO / LARS) &
Kernel methods (SVM path) & Control (explicit MPC) \\
\midrule
Homotopy parameter $\lambda$ & risk--return trade-off & penalty $\lambda$ &
cost $C$ & state $x$ (a \emph{vector}) \\[2pt]
Quadratic form $H$ (as operator) & covariance $\Sig$ & Gram matrix $X\T X$ &
kernel Gram $\kappa(X,X)$ & state--input cost \\[2pt]
Active set & assets held off their bounds (the free set) &
the support (active coefficients) & margin support vectors &
constraints active on a critical region \\[2pt]
Face mechanism & imposed by box bounds & induced by the $\ell_1$ subgradient &
induced by the margin & imposed by state/input constraints \\[2pt]
Breakpoint (event) & turning point / corner portfolio & knot / LARS event &
a point crossing the margin & critical-region boundary \\[2pt]
Event selection (ratio test) & next asset to enter or leave the free set &
next variable to enter or leave & next point to cross the margin &
next region boundary reached \\[2pt]
Solution path & efficient frontier & regularization path & regularization path &
explicit (piecewise-affine) control law \\
\bottomrule
\end{tabular}
\end{center}

The control column describes the \emph{multi-parameter} generalisation of
Remark~\ref{rem:vector}: the interval of breakpoints becomes a polyhedral partition into
critical regions and the homotopy becomes region traversal. The remaining entries are
its scalar slice.

\section{The four events, by dialect}

Chapter~\ref{ch:parametric} named four event families and Table~\ref{tab:events} read
two dialects. For reference, the full sorting:

\begin{center}
\footnotesize
\begin{tabular}{@{}llll@{}}
\toprule
Method & Events populated & Active set & Face mechanism \\
\midrule
Forward LARS               & $\mathrm{D}_1$
                           & monotone      & induced \\
LASSO homotopy             & $\mathrm{P}_1$, $\mathrm{D}_1$
                           & non-monotone  & induced \\
Critical Line Algorithm    & $\mathrm{P}_1$, $\mathrm{D}_1$
                           & non-monotone  & imposed (box) \\
CLA with coupling rows     & all four
                           & non-monotone  & imposed \\
Constrained LASSO          & all four
                           & non-monotone  & mixed \\
SVM path                   & $\mathrm{P}_1$, $\mathrm{D}_1$
                           & non-monotone  & induced (margin) \\
Explicit MPC               & all four, over a vector parameter
                           & non-monotone  & imposed \\
\bottomrule
\end{tabular}
\end{center}

The Critical Line Algorithm and the LASSO homotopy populate the same two families. The
difference between them is that the CLA's feasibility intervals are \emph{fixed} (a box)
while the LASSO's travel with $\lambda$ (the interval $[-\lambda,\lambda]$). That single
difference is all that separates an imposed active set from an induced one at the level
of the algorithm, and Theorem~\ref{thm:identity} is what makes the two interconvertible.
